\section{A specified boundary theory and normalized slit observable}
\label{sec:boundary}

\subsection{Boundary sectors and the chosen insertions}

Use the diagonal unitary $\su_2$ WZW theory, with central charge $c=3/2$,
maximally symmetric Cardy boundary conditions, and identity current gluing.
The integrable labels are $0,1/2,1$. Writing $V_j$ for the affine module
including descendants, the boundary sectors are
\begin{equation}
 \mathcal H_{ab}=\bigoplus_j N_{aj}^{\ b}V_j,
 \qquad \mathcal H_{00}=V_0,\quad
 \mathcal H_{0,1/2}=V_{1/2},\quad
 \mathcal H_{1/2,1/2}=V_0\oplus V_1.
 \label{eq:sectors}
\end{equation}
These are infinite-dimensional spaces. The boundary fusion input is the
standard Cardy construction \cite{Cardy,BPPZ}; it is distinct from the
two-dimensional invariant tensor space used below.

Let $D_t$ be the upper half-plane with a prescribed simple slit removed.
Its hydrodynamically normalized map satisfies
\begin{equation}
 \dot g_t(z)=\frac{2}{g_t(z)-u_t},\qquad g_0(z)=z,\qquad u_0=0.
 \label{eq:chordal}
\end{equation}
Assume $u$ is $C^1$ and work on an interval where the slit, marked prime
ends and ordered real spectators $0<x<y$ remain well defined and unswallowed.
The successive boundary labels in the uniformized half-plane are
\begin{equation}
 0\ \xrightarrow{u_t}\ \tfrac12
 \ \xrightarrow{g_t(x)}\ 0
 \ \xrightarrow{g_t(y)}\ \tfrac12
 \ \xrightarrow{\infty}\ 0.
 \label{eq:labels}
\end{equation}
Each arrow carries the unique spin-$1/2$ boundary-changing primary in the
corresponding sector. All four fields have weight $h=3/16$. They are
boundary fields, so there are no additional mirror insertions.

\subsection{The normalization is part of the observable}

Divide the four-field tensor amplitude by the two-field amplitude with the
same tip and infinity fields, whose invariant contraction is normalized to
one in the half-plane. Fix the remaining field normalization by the unit
leading vacuum OPE coefficient used in Section~\ref{sec:blocks}. Set
\begin{equation}
 a_t=g_t(x)-u_t,\quad b_t=g_t(y)-u_t,\quad
 J_x=g_t'(x),\quad J_y=g_t'(y).
\end{equation}
If $F$ is the normalized half-plane tensor amplitude, the slit observable is
\begin{equation}
 M_t(x,y)=J_x^hJ_y^h F(0,a_t,b_t,\infty).
 \label{eq:observable}
\end{equation}
Use the same uniformizing prime-end coordinate $g_t(z)-u_t$ and cutoff at
the tip in numerator and denominator. Their tip-primary factors then
cancel; no finite Euclidean derivative at the slit tip is invoked.
The infinity insertion is $\lim_{z\to\infty}z^{2h}\phi(z)$, and its
reciprocal-coordinate derivative is one by hydrodynamic normalization.
Only the two spectator Jacobians remain. The common domain/Weyl anomaly
factor cancels in the normalized ratio. There is no additional freely
adjustable scalar function of $t$.

This uses the normalized slit-correlator framework of
Alekseev--Bytsko--Izyurov \cite{ABI}. Here the driver is deterministic;
neither a Brownian driver nor an internal-group average, an It\^o term,
or a stochastic martingale theorem is imported.

\subsection{What the Chern--Simons interpretation supplies}

In holomorphic polarization on the doubled disk, the four marked fields
give a sphere with four spin-$1/2$ punctures. The $\su_2$ Chern--Simons
conformal-block space has the two allowed intermediate channels $0,1$.
A framed network preparing the first-pair vacuum channel selects the state
used below. The relation between Wilson endpoints and conformal-block state
spaces is established in the boundary formalism of
Alekseev--Barmaz--Mnev \cite{ABM}.

Fix ribbon framings and local coordinates throughout. This supplies a
Chern--Simons interpretation of the chiral data. It does not identify the
two-dimensional slit with a bulk Wilson trajectory, or identify the finite
channel space with either an affine boundary module or $L^2$ of an interval.
